\section{Deferral and Boundary}

Both layers are deferred. The downstream evidence specification explicitly says
not to implement RSTAT, RROLL, or the other downstream packages until their
access patterns and claim boundaries keep earlier packages from absorbing their
responsibilities \citep{civic-evidence-boundaries-2026}. That deferral is part
of the design, not a missing implementation note.

RSTAT should be built only after enough normalized RCOUNT and RHIST inputs exist
to make statistical analysis reproducible rather than anecdotal. A forensic
report over unstable or one-off inputs would encourage readers to treat a model
output as stronger than the evidence beneath it. RSTAT therefore waits for
stable package hashes, normalized time series, and claim-boundary language that
keeps a statistic from becoming certification.

RROLL has a different reason to wait. Aggregate eligibility universes are useful
only if the disclosure boundary is known before fixtures exist. The fixture
ladder should not teach future implementers to publish rare ballot styles,
small cells, or private voter-level rows. RROLL therefore waits for privacy and
disclosure thresholds, plus source-hash and caveat rules, before any positive or
negative fixture is written.

The operational boundary is simple. RSTAT can identify patterns that need
explanation. RROLL can declare aggregate eligibility universes. Neither
certifies an election. Certification belongs to the packages and legal actions
that bind count verification, audit evidence, and official signoff. Analysis and
eligibility universes may inform those conversations, but they do not replace
them.
